\documentclass[11pt]{article}

\usepackage[margin=1in]{geometry}
\usepackage{amsmath,amssymb,amsfonts}
\usepackage{booktabs}
\usepackage{enumitem}
\usepackage{graphicx}
\usepackage{hyperref}
\usepackage[numbers,sort&compress]{natbib}

\title{Article Plan: Reduced-Parameter Detection of Ring-Like Structures in Supermassive Black Hole VLBI Images}
\author{Draft outline}
\date{}

\begin{document}
\maketitle
\tableofcontents

\section{Working title and central claim}

\subsection{Possible titles}
\begin{enumerate}[leftmargin=*]
    \item Reduced-Parameter Detection of Ring-Like Morphologies in Horizon-Scale VLBI Images.
    \item Morphology-First Detection of Annular Structures in Reconstructed Black Hole VLBI Images.
    \item A Model-Constrained Computer-Vision Detector for Ring-Like Features in Event-Horizon-Scale Images.
    \item Statistical Cataloguing of Ring-Like Structures in GRMHD Black Hole Images.
\end{enumerate}

\subsection{Central claim}
The paper should present the method as a detection algorithm for annular morphology in VLBI images, not as a direct replacement for full physical inference. The method detects whether a reconstructed or simulated image contains a ring-like structure compatible with an analytical morphology family. The method searches a reduced parameter space, obtains candidates, and then fits the remaining nuisance parameters. The output is a detection decision, an interpretable set of ring parameters, uncertainty estimates, image-domain quality measures, and a catalog of candidate ring-like components.

\subsection{Claim boundaries}
The method should not claim direct detection of a photon ring unless the test data and observational resolution support that interpretation. The safer terms are ring-like structure, annular morphology, lensing-ring candidate, or thin ring-like component. The term photon ring should be used for the physical interpretation after additional tests against ray-traced simulations, visibility-domain residuals, theoretical subring predictions, and astrophysical alternatives.

\section{Abstract plan}

The abstract should contain six points. First, horizon-scale VLBI images of M87* and Sgr A* show robust annular emission, but extracting ring parameters depends on imaging assumptions, geometric model choices, and visibility-domain likelihoods. Second, photon rings are theoretically important because they provide a compact probe of strong-field lensing and Kerr geometry. Third, the proposed algorithm searches for annular morphology in a reduced parameter space and fits the remaining parameters only for high-evidence candidates. Fourth, the method is evaluated on synthetic VLBI-like images, null images, model-mismatch tests, and image reconstructions from simulated visibility data. Fifth, the method can be used on large GRMHD image libraries to build distributions of detected ring-like structures and to study how their parameters depend on black-hole mass, spin, inclination, emission model, and observing array. Sixth, the paper reports computer-vision detection metrics, VLBI residual metrics, parameter-estimation errors, calibration diagnostics, and runtime scaling.

\section{Introduction}

\subsection{Scientific context}
The Event Horizon Telescope has produced horizon-scale images of M87* and Sgr A*. The reported images are dominated by compact ring-like emission with a central brightness depression. These structures are consistent with strong gravitational lensing near a black hole and with emission from relativistic plasma near the event horizon \citep{EHTM87I2019,EHTM87IV2019,EHTSgrAI2022,EHTSgrAIII2022}. The 2018 M87* observations further support a persistent ring diameter across epochs, while the azimuthal brightness distribution changes with time \citep{EHTM87Persistent2024}.

\subsection{Need for morphology detection}
Current EHT analyses combine visibility-domain modeling, regularized maximum likelihood imaging, CLEAN imaging, Bayesian imaging, feature extraction, and comparison with GRMHD simulations \citep{EHTM87IV2019,EHTM87VI2019,EHTSgrAIV2022}. These methods are scientifically strong, but the detection of a specific annular morphology is not always separated from image reconstruction, astrophysical modeling, and parameter estimation. A dedicated morphology detector can provide an independent image-domain test of ring presence, ring geometry, and failure modes.

\subsection{Problem with overfitting and model dependence}
A detector trained or tuned only on one synthetic data set can obtain high precision and recall without being robust to different backgrounds, noise, reconstruction algorithms, or astrophysical morphologies. A full geometric or physical model fit can also be sensitive to the assumed source family, likelihood, priors, calibration treatment, and image regularization. The article should therefore state that the algorithm is not validated by one aggregate score. It is validated by controlled tests across signal-to-noise ratio, angular resolution, beam convolution, uv coverage, ring width, asymmetry, background emission, dynamical variability, and model mismatch.

\subsection{Contributions}
\begin{enumerate}[leftmargin=*]
    \item Define ring detection in reconstructed SMBH VLBI images as a hypothesis test between a null image and an annular morphology model.
    \item Introduce a reduced-parameter search over position, scale, and optional ellipticity, followed by conditional fitting of width, amplitude, asymmetry, and local background parameters.
    \item Provide a reproducible algorithm with candidate generation, robust model fitting, duplicate rejection, multiple-testing control, and calibrated detection scores.
    \item Evaluate the method using computer-vision metrics and VLBI-specific diagnostics.
    \item Compare against annular Hough detection, randomized Hough detection, template matching, EHT-style image feature extraction, and geometric model fitting where feasible.
    \item Use the detector as a statistical tool for GRMHD image libraries, where detected ring parameters become summary statistics for black-hole and accretion-model inference.
\end{enumerate}

\section{Literature overview}

\subsection{Black hole shadows and ring-like appearance}
The theoretical basis for ring-like black hole images comes from relativistic ray tracing and gravitational lensing near compact objects. Early calculations showed the optical appearance of a black hole with an accretion disk \citep{Luminet1979}. Submillimeter VLBI was later proposed as a way to resolve the shadow of Sgr A*, with the expected angular scale close to the achievable resolution of long-baseline interferometry \citep{Falcke2000}. The EHT results for M87* and Sgr A* confirmed that horizon-scale data can recover compact annular emission with a central depression \citep{EHTM87I2019,EHTM87IV2019,EHTM87VI2019,EHTSgrAI2022,EHTSgrAIII2022,EHTSgrAIV2022}.

\subsection{Photon rings, lensing rings, and shadows}
The terminology must be precise. The observed bright annulus in a finite-resolution VLBI image is not automatically the mathematical photon ring. The photon ring is associated with highly bent photon trajectories near the critical curve. The observed ring can include direct emission, lensing-ring emission, beam effects, and reconstruction effects. \citet{GrallaHolzWald2019} clarify the distinction between black hole shadows, photon rings, and lensing rings. \citet{Johnson2020} show that photon-ring subrings can produce universal interferometric signatures on long baselines. \citet{GrallaLupsascaMarrone2020} argue that the photon-ring shape can be a precise test of strong-field general relativity. Later work studies finite-width effects, observing requirements, thick-disk emission, and departures from the universal photon-ring regime \citep{Wielgus2022,Vincent2022,Paugnat2022,Jia2024}.

\subsection{Photon rings as probes of mass and spin}
The article should include a short review of how photon-ring observables connect to black-hole parameters. The angular diameter of a strong-lensing ring is controlled mainly by the gravitational radius divided by distance, with weaker dependence on spin and viewing geometry. Shape, asymmetry, relative subring size, polarimetric signatures, and long-baseline interferometric signatures can carry additional information about spin and inclination \citep{GrallaLupsascaMarrone2020,Broderick2022,Palumbo2023}. This review should also state the limitation. Ring parameters measured from finite-resolution reconstructed images are not identical to theoretical photon-ring parameters, so physical inference must be calibrated with ray-traced GRMHD simulations.

\subsection{EHT imaging and feature extraction}
EHT image reconstruction is an ill-posed inverse problem because VLBI samples sparse Fourier components of the sky brightness. Classical radio interferometric imaging includes CLEAN deconvolution and maximum entropy methods \citep{Hogbom1974,Cornwell1985,Sault1990,Cornwell2009,Thompson2017}. EHT analyses use CLEAN, regularized maximum likelihood methods, closure quantities, Bayesian methods, and imaging surveys over reconstruction assumptions \citep{Chael2018,EHTM87IV2019,EHTSgrAIII2022}. Computational imaging methods for VLBI provide an additional context for image-domain morphology measurements \citep{Bouman2016}. Recent deep-learning reconstruction with closure invariants should be cited as a modern reconstruction baseline or future comparison, not as the main baseline unless it is implemented in the experiments \citep{Lai2025}. Because non-parametric images do not directly return ring parameters, feature-extraction tools such as REx and VIDA are used to measure diameter, width, orientation, asymmetry, and ellipticity from reconstructed images \citep{TiedeVIDA2022,TiedeEllipticity2022}. This provides the closest image-domain comparison for the proposed detector.

\subsection{Visibility-domain geometric modeling}
EHT analyses also fit geometric and GRMHD models directly to interferometric data. For M87*, geometric crescent models and GRMHD models were explored with MCMC, nested sampling, and genetic algorithms using visibility amplitudes, closure phases, and log closure amplitudes \citep{KamruddinDexter2013,EHTM87VI2019}. For Sgr A*, geometric ring models were fitted using strategies that account for source variability, and the ring diameter was measured by both geometric modeling and image-domain feature extraction \citep{EHTSgrAIV2022}. This class of methods is a strong baseline, but it is model-dependent and computationally expensive for large model families.

\subsection{Photon-ring measurements and hybrid imaging}
Hybrid imaging models a flexible image component together with a geometric ring component. It is directly relevant to photon-ring searches because it tries to separate diffuse emission from a thin ring-like component. \citet{BroderickPhotonRing2022} report a thin secondary-image component in M87* using simultaneous modeling and imaging. \citet{TiedeNgEHT2022} caution that hybrid imaging can produce false positives for photon rings and that fitted ring parameters may not directly correspond to physical photon-ring properties. This motivates explicit false-positive tests, null experiments, and false-discovery-rate control in the proposed work \citep{BenjaminiHochberg1995}.

\subsection{Future arrays and space VLBI}
The proposed detector is relevant for next-generation and space-based arrays because higher angular resolution increases the number of ring-like substructures that can be tested. ngEHT science goals include sharper imaging, time-domain studies, and photon-ring science \citep{Johnson2023,LupsascaBlackHoleExplorer2024}. Space-VLBI concepts such as Millimetron, the Event Horizon Imager, PECMEO, Capella, and related subterahertz developments provide additional motivation for algorithms that can process many reconstructions and compare ring statistics across uv coverage and observing geometry \citep{Novikov2021,Rudnitskiy2021,Rudnitskiy2023,Martin2019,Kudriashov2019,Kudriashov2021,Roelofs2019,Trippe2023,Likhachev2024,Shlentsova2024}.

\subsection{Computer-vision background}
The proposed method is related to classical pattern recognition and computer vision. Hough-transform methods detect low-dimensional geometric structures through voting in parameter space \citep{DudaHart1972}. Circle and ellipse detection literature is directly relevant because the target morphology is annular rather than a semantic object category. The algorithm section should cite the comparative circle-Hough review and randomized Hough transform when motivating the reduced search and candidate generation steps \citep{YUEN199071,Xu1990}. The algorithm section should cite probabilistic Hough-transform work when discussing speed, sparse voting, and statistical candidate selection \citep{KALVIAINEN1995239}. RANSAC detects models by hypothesis generation, scoring, and refitting \citep{FischlerBolles1981}. Persistent homology can be cited as an optional topology-based diagnostic for ring persistence under thresholding, but not as a required component unless implemented \citep{Bauer2014}. Modern detection evaluation reports precision, recall, average precision, localization accuracy, and robustness under distribution shift \citep{Everingham2010,LinCOCO2014}. The paper should use these metrics so that the algorithm is legible to both VLBI and computer-vision audiences.

\section{Problem formulation}

\subsection{Data model}
Let $I(x)$ be a reconstructed total-intensity VLBI image on a two-dimensional angular grid, where $x=(x_1,x_2)$ is measured in microarcseconds. Let $P$ be the effective restoring beam or point-spread operator associated with the reconstruction. The observed image is modeled as
\begin{equation}
    I(x) = P \ast \left[B(x) + A(x;\theta,\eta)\right] + \epsilon(x),
\end{equation}
where $B(x)$ is local background emission, $A(x;\theta,\eta)$ is the analytical ring morphology, and $\epsilon(x)$ includes reconstruction noise, thermal noise propagated to image space, residual calibration artifacts, and model mismatch.

\subsection{Ring model}
A practical analytical model is an elliptical annulus with angular brightness modulation:
\begin{equation}
    A(x;\theta,\eta) = a(\varphi) \exp\left[-\frac{(r_e(x;c,q,\phi)-R)^2}{2w^2}\right],
\end{equation}
where $c=(c_x,c_y)$ is the center, $R$ is the mean radius, $q$ is the axis ratio, $\phi$ is the position angle, $w$ is the ring width, and $a(\varphi)$ is the azimuthal brightness profile. A compact choice is
\begin{equation}
    a(\varphi)=a_0\left[1+\sum_{k=1}^{K}\alpha_k\cos(k\varphi)+\beta_k\sin(k\varphi)\right].
\end{equation}
The model should be convolved with the effective beam before comparison with the reconstructed image. The text should also state whether the fitted structure is a filled annulus, a crescent-like annulus, a thin ring component, or a nested set of annular components.

\subsection{Multiple-ring extension}
For GRMHD images and higher-resolution arrays, the image can contain several ring-like components. The model can be extended to
\begin{equation}
    I(x)=P\ast\left[B(x)+\sum_{j=1}^{J} A_j(x;\theta_j,\eta_j)\right]+\epsilon(x).
\end{equation}
The integer $J$ should not be fixed by the algorithm unless the experiment requires it. The detector should first return a candidate list and then merge, rank, and label candidates. Candidate labels should be descriptive: primary annulus, secondary annulus, inner ring-like component, outer ring-like component, or photon-ring candidate. Physical labels should be assigned only after comparison with ray-traced ground truth or theoretical expectations.

\subsection{Parameter split}
The full parameter vector is divided into search parameters and fitted nuisance parameters:
\begin{equation}
    t=(\theta,\eta), \qquad \theta=(c_x,c_y,R,q,\phi), \qquad \eta=(w,a_0,\alpha_k,\beta_k,B,\sigma).
\end{equation}
A simpler first implementation can search only $\theta=(c_x,c_y,R)$ and fit $q,\phi,w,a_0,\alpha_k,\beta_k$ afterward. The article should justify the chosen split by computational cost, identifiability, robustness, and expected VLBI resolution.

\subsection{Detection hypotheses}
Define a null hypothesis and an annular morphology hypothesis:
\begin{align}
    H_0 &: I(x)=P\ast B(x)+\epsilon(x), \\
    H_1 &: I(x)=P\ast[B(x)+A(x;\theta,\eta)]+\epsilon(x).
\end{align}
The detection problem is to decide whether $H_1$ is supported and to estimate $\theta$ and $\eta$ when it is supported. For multiple candidates, the problem is also to control false discoveries across the searched parameter space.

\subsection{Reduced evidence function}
The core score is a profile evidence over the reduced parameter space:
\begin{equation}
    S(\theta)=\max_{\eta}\, Q\left(I, P\ast [B+A(\cdot;\theta,\eta)]\right),
\end{equation}
where $Q$ may be a negative robust residual, a likelihood, a normalized annular contrast, or a calibrated probability. In implementation, $S(\theta)$ can be approximated by a fast annular matched score before the full fit. This connects the method to profile likelihood, nuisance-parameter elimination, Hough-style voting, and RANSAC-style proposal generation.

\section{Algorithm}

\subsection{Inputs and outputs}
Inputs are a reconstructed VLBI image, pixel scale, beam description, optional image uncertainty map, candidate parameter ranges, and thresholds. Optional inputs are visibility-domain residuals, reconstruction metadata, and GRMHD ground-truth labels. Outputs are detected ring candidates, fitted parameters, confidence scores, annular masks, residual images, candidate catalogs, and diagnostic statistics.

\subsection{Step-by-step procedure}
\begin{enumerate}[leftmargin=*]
    \item Preprocess the image. Normalize flux. Apply optional beam-aware smoothing. Mask unreliable image regions. Estimate the local background scale.
    \item Compute an annular evidence map $S(c_x,c_y,R)$ or $S(c_x,c_y,R,q,\phi)$ using normalized contrast between an annular band and the central depression plus exterior background. This part should cite circle Hough, randomized Hough, and probabilistic Hough methods because the detector also uses a reduced parameter-space search \citep{YUEN199071,Xu1990,KALVIAINEN1995239}.
    \item Detect local maxima in $S$. Apply minimum separation in center and radius. Keep the top candidates or all candidates above a threshold. Correct for multiple candidate tests using validation-calibrated thresholds or false-discovery-rate control \citep{BenjaminiHochberg1995}.
    \item For each candidate, fit the full beam-convolved annular model using robust least squares, profile likelihood, or a Bayesian local fit. Fit the local background simultaneously.
    \item Compute quality statistics. Use annular contrast, central-depression significance, residual energy, parameter uncertainty, ring-coverage fraction, and optional topological persistence under threshold changes \citep{Bauer2014}.
    \item Reject candidates that fail physical or statistical constraints. Reject candidates with excessive width, low central-depression significance, poor residuals, unstable fitted parameters, or insufficient angular coverage.
    \item Merge duplicate detections with non-maximum suppression in parameter space. Prefer the candidate with the higher calibrated score.
    \item Return the final detection list, fitted parameter covariance, annular mask, residual image, and diagnostic plots.
\end{enumerate}

\subsection{Decision score}
The final score should combine fast evidence and fit quality:
\begin{equation}
    C = \lambda_1 \tilde{S} + \lambda_2 Z_{\rm dep} + \lambda_3 Z_{\rm ann} - \lambda_4 \chi^2_{\rm red} - \lambda_5 U_t - \lambda_6 F_{\rm cov}^{-1},
\end{equation}
where $Z_{\rm dep}$ is central-depression significance, $Z_{\rm ann}$ is annular-brightness significance, $\chi^2_{\rm red}$ is reduced residual error, $U_t$ is normalized parameter uncertainty, and $F_{\rm cov}$ is the angular coverage fraction of the detected annulus. The weights should be fixed on a validation set and not changed on the final test set.

\subsection{Photon-ring interpretation gate}
A detection should be called a photon-ring candidate only if it passes additional tests. The fitted diameter and width must be compatible with ray-traced black-hole models after beam convolution. The component must survive uv-coverage and reconstruction tests. The residual visibilities or closure quantities must not degrade. The same component should be stable under reconstruction choices. Null and model-mismatch experiments must show controlled false-positive rates. If the algorithm finds several nested ring-like components, the photon-ring label should be assigned only when simulations or theoretical subring predictions support that interpretation \citep{Johnson2020,Wielgus2022,Vincent2022,Jia2024}.

\section{Using the detector for statistical insight from GRMHD simulations}

\subsection{Purpose}
This section should be placed before the experimental validation because it defines a scientific use case beyond single-image detection. The algorithm can be used to convert a large library of GRMHD and ray-traced images into a catalog of ring-like morphology parameters. This makes it possible to study which ring-like structures are common, which are rare, and how their measured parameters depend on black-hole and accretion parameters.

\subsection{Simulation program}
Generate a large synthetic library from GRMHD simulations and ray tracing. Vary black-hole mass $M$, spin $a_*$, inclination $i$, distance $D$, electron-heating prescription, magnetic state, observing frequency, optical depth, time snapshot, and scattering model when relevant. For each sky image, also generate beam-convolved images and, when possible, reconstructed images from synthetic VLBI observations. The detector should be run on all image products, not only on selected visually clear cases.

\subsection{Detected-ring catalog}
For each simulated image, find all statistically supported ring-like structures. Store center, diameter, width, ellipticity, position angle, azimuthal asymmetry, central-depression score, ring coverage, fit residual, confidence score, and uncertainty. If ray-tracing metadata are available, label detections by emission origin or lensing order. If no such metadata are available, use neutral labels and avoid calling a component a photon ring only from image morphology.

\subsection{Inference target}
The detected-ring catalog can be used to estimate integral black-hole parameters statistically. Define a summary vector
\begin{equation}
    z = (d, w, q, \phi, A_1, A_2, Z_{\rm dep}, Z_{\rm ann}, C),
\end{equation}
where $d=2R$ is the diameter and $A_k$ are asymmetry coefficients. Use the simulation library to estimate
\begin{equation}
    p(z \mid M,a_*,i,D,\mathcal{E}),
\end{equation}
where $\mathcal{E}$ denotes the emission and accretion model. For observed data, compare the measured feature vector $z_{\rm obs}$ with this simulation-calibrated distribution. The inference should be written as statistical matching to simulated observables, not as a direct formula from a single ring diameter to spin.

\subsection{Why this is useful}
The approach is useful for four reasons. First, it produces a compact and interpretable summary of thousands of GRMHD images. Second, it separates morphology statistics from the full high-dimensional image comparison problem. Third, it allows sensitivity studies: one can test which detected features change with mass, spin, inclination, emission model, and uv coverage. Fourth, it gives an empirical prior for later visibility-domain fitting by identifying plausible ring diameters, widths, asymmetries, and candidate substructures. The method can support spin and mass inference when the simulation library shows that the measured features are informative and when degeneracies with distance, inclination, and emission physics are included.

\subsection{Validation inside the GRMHD study}
The GRMHD section should include its own validation. Use simulations with known ray-tracing labels to measure completeness and purity for direct emission rings, lensing rings, and possible photon-ring components. Report feature sensitivity to spin and mass with confidence intervals. Test whether inferred parameters are unbiased on held-out simulations. Report posterior coverage if the catalog is used for inference. Include a null test where the detector is run on GRMHD images that do not contain a resolvable thin ring component.

\subsection{Figures and tables for this section}
Required figures are a grid of GRMHD images with detected rings, distributions of detected diameter and width across simulations, feature trends versus spin and inclination, feature trends versus mass-to-distance ratio, confusion matrices for ring labels when labels are available, and posterior or likelihood contours for $M$ and $a_*$. Required tables are the simulation parameter grid, the detected-ring catalog schema, feature importance or sensitivity ranking, held-out inference errors, and false-positive rates for unresolved or non-ring cases.

\section{Why the method can improve on current techniques}

\subsection{Careful wording}
The paper should not state that the method is universally better than EHT maximum-likelihood or Bayesian modeling. A defensible claim is narrower. The method can be better as a detection-stage and cataloguing tool when the goal is to find annular morphology with limited assumptions about the emission model.

\subsection{Arguments to test}
\begin{enumerate}[leftmargin=*]
    \item Reduced dimensionality. The algorithm searches a low-dimensional space first and fits nuisance parameters only for plausible candidates. This can reduce computation relative to blind full-parameter fitting.
    \item Detection before interpretation. The method separates the question of whether an annular morphology exists from the question of which physical accretion model produced it.
    \item Explicit null control. The method reports false positives per image and false positives per null data set. This is necessary because hybrid ring fitting can produce false positives in some regimes \citep{TiedeNgEHT2022}.
    \item Image-domain interpretability. The output is a ring mask, fitted center, diameter, width, ellipticity, asymmetry, and residual map. These are easy to inspect and compare across reconstructions.
    \item Robustness reporting. The method reports performance as a function of SNR, beam size, uv coverage, background complexity, reconstruction pipeline, and model mismatch.
    \item Compatibility with EHT methods. The algorithm can be used after CLEAN, RML, Bayesian imaging, deep-learning reconstruction, or synthetic reconstructions. It can also produce candidate initializations for visibility-domain model fitting \citep{Hogbom1974,Chael2018,Bouman2016,Lai2025}.
    \item Simulation-scale use. The algorithm can be applied to large GRMHD libraries where full visibility-domain fitting for every candidate and every snapshot is expensive. This is a practical advantage for population-level morphology studies.
\end{enumerate}

\subsection{Comparison targets}
The article should compare against at least five baselines. The first is annular Hough or ellipse Hough detection. The second is randomized or probabilistic Hough detection. The third is matched filtering with fixed templates. The fourth is RML image reconstruction followed by REx or VIDA feature extraction. The fifth is geometric crescent fitting to visibility data where simulated visibilities are available. If direct visibility-domain fitting is not implemented, the article should state that the comparison is image-domain only.

\section{Experiments}

\subsection{Synthetic image-domain experiments}
Generate images with known annular ground truth:
\begin{equation}
    I(x)=P\ast\left[B(x)+A(x;t) + D(x)\right]+\epsilon(x),
\end{equation}
where $D(x)$ represents non-ring distractors such as compact knots, crescent arcs, jet-like extensions, and smooth asymmetric emission. Randomize center, radius, width, ellipticity, position angle, brightness asymmetry, flux scale, background gradient, beam, noise, and image size.

\subsection{Null experiments}
Generate images with no true annulus. Use smooth disks, crescents without central depression, compact Gaussian components, jets, arcs, turbulent backgrounds, and reconstruction-like artifacts. Report the false-positive rate, false detections per image, and false-discovery-rate behavior. This experiment is essential for credibility.

\subsection{Model-mismatch experiments}
Test images generated from models that differ from the detector model. Use thick crescent rings, partial rings, non-Gaussian radial profiles, elliptical rings with angularly varying width, nested rings, and GRMHD-like brightness maps. This tests whether the algorithm detects general annular morphology rather than only its exact formula.

\subsection{VLBI forward-simulation experiments}
Use ray-traced or GRMHD images as sky truth. Sample synthetic visibilities using EHT-like, ngEHT-like, and space-VLBI-like uv coverage when possible. Add thermal noise, station gain errors, atmospheric phase corruption, and for Sgr A* include scattering and variability scenarios when relevant. Reconstruct the images using at least two pipelines, such as CLEAN and RML. Run the detector on the reconstructions and compare recovered parameters to the known truth after beam convolution \citep{Roelofs2019,Johnson2023,Kudriashov2021,Trippe2023}.

\subsection{GRMHD population experiments}
Run the detector on a large library of GRMHD snapshots and ray-traced images. Report how often the detector finds zero, one, or several ring-like structures. Report distributions of diameter, width, ellipticity, and asymmetry as functions of $M$, $a_*$, inclination, observing frequency, and emission model. Use held-out simulations to test whether the detected morphology features can recover mass-to-distance ratio and constrain spin. This experiment should be separate from the basic detection benchmark because it tests scientific utility rather than only detection accuracy.

\subsection{Real-data demonstration}
Apply the detector to published or reproducible reconstructions of M87* and Sgr A* only as a demonstration unless raw data and complete calibration are included. Report the fitted ring parameters and residual maps. Do not use real-data agreement alone as proof of correctness.

\section{Metrics}

\subsection{Computer-vision detection metrics}
Report object-level precision, recall, F1 score, average precision, mean average precision across match thresholds, ROC-AUC if negatives are well defined, and false positives per image. Use a detection match criterion based on center error, radius error, and annular-mask overlap. A suggested criterion is
\begin{equation}
    \frac{\lVert \hat{c}-c\rVert}{R} < \tau_c, \qquad \frac{|\hat{R}-R|}{R}<\tau_R, \qquad {\rm IoU}(M,\hat{M})>\tau_{\rm IoU}.
\end{equation}
Report PR curves and FROC curves because false detections per image are important in scientific imaging.

\subsection{Localization and morphology metrics}
Report center error, diameter error, width error, position-angle error, ellipticity error, annular-mask IoU, Dice coefficient, and Hausdorff distance for the outer and inner contours when masks are available. Report median, interquartile range, and 95th percentile, not only means.

\subsection{VLBI-specific metrics}
When visibilities are available, report visibility-amplitude residuals, closure-phase residuals, log-closure-amplitude residuals, reduced chi-square, and residuals as a function of baseline length. Also report whether adding the fitted ring component improves or worsens the fit to interferometric observables.

\subsection{Multiple-testing and calibration metrics}
Report empirical false discovery rate, reliability diagrams, expected calibration error, Brier score, and coverage of confidence intervals. A probability of 0.8 should correspond approximately to an 80 percent empirical detection frequency under the test distribution.

\subsection{GRMHD inference metrics}
For simulation-based inference, report mass error, spin error, inclination error, posterior coverage, rank histograms, calibration curves, and sensitivity of features to nuisance parameters. If the method only gives a feature catalog and not a posterior, report correlation, mutual information estimates, or regression error between detected features and simulation parameters.

\subsection{Runtime metrics}
Report wall-clock time per image, runtime versus image size, runtime versus number of radius samples, runtime versus number of candidates, runtime versus number of GRMHD snapshots, and memory usage. Compare reduced search plus fitting against full grid search where feasible.

\section{Required figures}

\begin{enumerate}[leftmargin=*]
    \item Conceptual figure. Show reconstructed VLBI image, analytical ring model, reduced search space, candidate peak, and final fitted ring.
    \item Literature-positioning figure. Show the analysis chain from VLBI observables to imaging, feature extraction, geometric fitting, hybrid imaging, GRMHD simulation comparison, and the proposed morphology detector.
    \item Algorithm pipeline. Show preprocessing, evidence map, candidate extraction, full fit, rejection, and output.
    \item Synthetic data examples. Show positive examples, null examples, distractor examples, and model-mismatch examples.
    \item Evidence maps. Show $S(c_x,c_y,R)$ slices and peaks for true positives and false positives.
    \item Detection overlays. Show fitted annuli on easy, medium, hard, and failed cases.
    \item Precision-recall curve. Include confidence bands from bootstrap resampling.
    \item FROC curve. Plot recall versus false positives per image.
    \item Metrics versus SNR. Plot precision, recall, F1, AP, and mAP against noise level.
    \item Metrics versus resolution. Plot diameter error and recall against beam FWHM or maximum baseline.
    \item Metrics versus ring width and asymmetry. Show where the method fails.
    \item Metrics versus uv coverage. Compare EHT-like, ngEHT-like, and space-VLBI-like simulations.
    \item Parameter error distributions. Use boxplots for center, diameter, width, ellipticity, and orientation.
    \item Baseline comparison. Show proposed detector, Hough, randomized Hough, template matching, REx or VIDA, and geometric fitting.
    \item Ablation plot. Remove central-depression score, full fit, beam convolution, background fit, duplicate rejection, and multiple-testing correction.
    \item Residual plots. Show image residuals and visibility-domain residuals if available.
    \item GRMHD catalog figure. Show detected rings across a grid of spin, inclination, and emission-model settings.
    \item GRMHD statistics figure. Plot detected diameter, width, ellipticity, and asymmetry versus $M/D$, spin, and inclination.
    \item GRMHD inference figure. Show likelihood or posterior contours for $M$ and $a_*$ from detected ring features on held-out simulations.
    \item Failure cases. Show representative false positives and false negatives with explanation.
\end{enumerate}

\section{Required tables}

\begin{enumerate}[leftmargin=*]
    \item Parameter table. List $c_x,c_y,R,q,\phi,w,a_0,\alpha_k,\beta_k$ with units, ranges, and search-or-fit status.
    \item Synthetic generator table. List image sizes, beam sizes, SNR ranges, flux ranges, background models, noise models, and number of images.
    \item Main detection table. Precision, recall, F1, AP, mAP, ROC-AUC, FROC at fixed false-positive rates, empirical FDR, and confidence intervals.
    \item Parameter accuracy table. Center, diameter, width, ellipticity, orientation, and asymmetry errors.
    \item Robustness table. Metrics grouped by SNR, beam FWHM, uv coverage, ring width, and background class.
    \item Baseline table. Proposed method versus Hough, randomized Hough, probabilistic Hough, template matching, REx or VIDA, and geometric fitting.
    \item Ablation table. Full method versus removed components.
    \item GRMHD simulation table. List the grid of $M$, $a_*$, inclination, emission model, frequency, scattering, and number of snapshots.
    \item GRMHD catalog table. Define each stored feature, uncertainty, quality flag, and physical label if available.
    \item GRMHD inference table. Report held-out errors and confidence intervals for $M$, $a_*$, inclination, and mass-to-distance ratio.
    \item Runtime table. Reduced search, full fit, total time, and scaling.
\end{enumerate}

\section{Statistical protocol}

Use separate development, validation, and final test sets. Tune thresholds only on the development and validation sets. Freeze all thresholds before the final test. Use at least five independent random seeds for the synthetic generator. Report bootstrap confidence intervals over images and over observing simulations. Include null-only tests. Include out-of-generator tests where the data generator differs from the detector model. For GRMHD studies, split simulations by physical parameters or by simulation run so that held-out tests are not just adjacent snapshots from the same model.

\section{Discussion plan}

The discussion should state what the method proves and what it does not prove. It can prove that an image contains a statistically significant annular morphology under the tested conditions. It cannot alone prove that the detected component is the photon ring. That interpretation requires consistency with theoretical photon-ring predictions, stability under reconstruction, and visibility-domain support. The method is best positioned as an independent morphology detector, a quality-control tool for reconstructions, a candidate generator for later physical modeling, and a statistical feature extractor for GRMHD simulation libraries.

\section{Conclusion plan}

The conclusion should state that the paper introduces a reduced-parameter detector for ring-like structures in SMBH VLBI images. The method is designed to be interpretable, computationally efficient, and testable by both computer-vision and VLBI standards. The final claim should be based on controlled experiments: detection accuracy, false-positive control, parameter recovery, robustness to VLBI effects, comparison to existing feature-extraction and model-fitting methods, and usefulness for GRMHD population studies.

\section{Supplementary material plan}

Include implementation details, hyperparameters, pseudocode, random seeds, extra PR and FROC curves, full synthetic generator specification, all ablations, additional failure cases, GRMHD catalog schema, and simulation split details. Include code for reproducing every figure and table.

\bibliographystyle{plainnat}
\bibliography{ring_detector_refs_updated}

\end{document}
